\documentclass[11pt]{article}

\usepackage[utf8]{inputenc}
\usepackage[T1]{fontenc}
\usepackage{lmodern}
\usepackage{geometry}
\usepackage{amsmath,amssymb,amsfonts,amsthm,mathtools}
\usepackage{bm}
\usepackage{enumitem}
\usepackage{microtype}
\usepackage{hyperref}
\usepackage{titlesec}
\usepackage{booktabs}
\usepackage{array}
\usepackage{setspace}
\usepackage{mathrsfs}
\usepackage{physics}

\geometry{margin=1in}
\hypersetup{colorlinks=true, linkcolor=black, urlcolor=blue, citecolor=black}
\setlist[enumerate]{topsep=0.4em,itemsep=0.35em}
\setlist[itemize]{topsep=0.3em,itemsep=0.25em}
\titleformat{\section}{\large\bfseries}{\thesection.}{0.5em}{}
\titleformat{\subsection}{\normalsize\bfseries}{\thesubsection.}{0.5em}{}
\setstretch{1.06}

\newcommand{\Aether}{\AE{}ther}
\newcommand{\Aflow}{\AE{}ther-flow}
\newcommand{\TheoryName}{\Aflow{} Interpretation of Relativity}
\newcommand{\TheoryProgram}{\Aether{} / \Aflow{} framework}

\newcommand{\CompletedMetric}{g^{\sharp}}

\newcommand{\Car}{\mathrm{car}}
\newcommand{\Cur}{\mathrm{cur}}
\newcommand{\Hom}{\mathrm{hom}}

\newcommand{\ForSpace}[1]{\mathcal I_{#1}^{\mathrm{for}}}
\newcommand{\PrimShell}{\mathcal S_{\mathrm{prim}}}

\newcommand{\Reduction}[1]{\mathcal R_{#1}}
\newcommand{\CurrentCarrier}[1]{\mathfrak C_{#1}^{\mathrm{cur}}}
\newcommand{\EqCarrier}[1]{\mathfrak C_{#1}^{\mathrm{eq}}}
\newcommand{\CurrentExtract}[1]{\mathscr E_{#1}^{\mathrm{cur}}}
\newcommand{\EqExtract}[1]{\mathscr E_{#1}^{\mathrm{eq}}}
\newcommand{\PrimCurrent}[1]{\mathcal J_{#1}^{\mathrm{prim}}}
\newcommand{\PrimTheta}[1]{\Theta_{#1}^{\mathrm{prim}}}

\newcommand{\MetricOut}{\mathscr G_{\mathrm{out}}}
\newcommand{\MatterOut}{\mathscr T_{\mathrm{out}}}

\newcommand{\VisibleAffine}[1]{\widehat X_{#1}}
\newcommand{\DataSpace}[1]{\mathcal Y_{#1}^{\mathrm{obs}}}
\newcommand{\AmbientTarget}[1]{E_{#1}^{\mathrm{amb}}}

\newcommand{\Kernel}[1]{K_{#1}^{0}}
\newcommand{\StateComp}[1]{P_{#1}^{\mathrm{co}}}

\newcommand{\StatSpace}[1]{\mathcal U_{#1}^{\mathrm{st}}}
\newcommand{\StatBody}[1]{\mathcal B_{#1}^{\mathrm{st}}}
\newcommand{\StatData}[1]{\Lambda_{#1}^{\mathrm{st,dat}}}
\newcommand{\StatShell}[1]{\Lambda_{#1}^{\mathrm{st,sh}}}
\newcommand{\StatPotential}[1]{\Phi_{#1}^{\mathrm{st}}}
\newcommand{\StatZero}[1]{V_{#1}^{\mathrm{st},0}}
\newcommand{\StatHidden}[1]{H_{#1}^{\mathrm{st}}}
\newcommand{\StatHiddenBody}[1]{D_{#1}^{\mathrm{st}}}
\newcommand{\StatProj}[1]{\Pi_{#1}^{\mathrm{st,hid}}}
\newcommand{\StatHess}[1]{\mathfrak H_{#1}^{\mathrm{st}}}
\newcommand{\StatSym}[1]{\mathbb T_{#1}}
\newcommand{\StatTime}[1]{\vartheta_{#1}^{\mathrm{st}}}
\newcommand{\StatEmbed}[1]{\zeta_{#1}^{\mathrm{st}}}

\newcommand{\RedSpace}[1]{\mathcal U_{#1}^{\mathrm{red}}}
\newcommand{\RedBody}[1]{\mathcal B_{#1}^{\mathrm{red}}}
\newcommand{\RedData}[1]{\Lambda_{#1}^{\mathrm{red,dat}}}
\newcommand{\RedShell}[1]{\Lambda_{#1}^{\mathrm{red,sh}}}
\newcommand{\RedTarget}[1]{F_{#1}^{\mathrm{red}}}
\newcommand{\RedMap}[1]{\Gamma_{#1}^{\mathrm{red}}}
\newcommand{\RedZero}[1]{V_{#1}^{\mathrm{red},0}}
\newcommand{\RedHidden}[1]{H_{#1}^{\mathrm{red}}}
\newcommand{\RedHiddenBody}[1]{D_{#1}^{\mathrm{red}}}
\newcommand{\RedProj}[1]{\Pi_{#1}^{\mathrm{red,hid}}}
\newcommand{\RedLin}[1]{\mathcal M_{#1}^{\mathrm{red}}}
\newcommand{\RedSym}[1]{\mathbb R_{#1}}
\newcommand{\RedTime}[1]{\vartheta_{#1}^{\mathrm{red}}}
\newcommand{\RedEmbed}[1]{\zeta_{#1}^{\mathrm{red}}}

\newtheorem{definition}{Definition}
\newtheorem{lemma}{Lemma}
\newtheorem{proposition}{Proposition}
\newtheorem{remark}{Remark}
\newtheorem{corollary}{Corollary}
\newtheorem{theorem}{Theorem}

\title{\textbf{The \TheoryName{}}\\[0.4em]
\large Primitive-Reservoir Observer-Reduction\\
\large Zero-Hidden Stationary Hessian Law\\
\large Necessity / Uniqueness from Zero-Hidden Stationary Reduction Law}
\author{}
\date{}

\begin{document}

\maketitle

\begin{abstract}
The current benchmark-facing theorem on the active primitive-reservoir observer-reduction line already sharpened the burden once more: the zero-hidden spectral generator law is no longer merely available, because it is derived canonically from a zero-hidden stationary Hessian law and is necessary there up to the obvious zero/hidden spectral identifications. The remaining honest task is therefore no longer to justify the spectral-law layer. It is to justify why the stationary-Hessian law itself should hold.

The present note takes the smallest honest step now available below that stationary-Hessian layer without claiming final microphysical closure. Rather than postulate a nonnegative stationary potential and its common zero-locus Hessian directly, it collapses that package into a smaller zero-hidden stationary reduction law. For each sector this deeper package consists of a compact convex reduction-state body, affine observer-data and shell-response readouts, a Euclidean stationary-reduction target, a smooth reduction map with an exact zero-hidden reduction locus, exact body splitting over zero-hidden and hidden directions, a common zero-locus linearization that vanishes on zero directions and is coercively positive on hidden directions, symmetry and time-reversal actions preserving the full reduction law, an exact zero-response shell realization on the zero locus, fixed-data shell solvability on that same locus, the inert zero-hidden kernel and projector, and a reduction coercivity inequality on data-invisible variations.

From that stationary reduction law one derives canonically every constituent of the stationary-Hessian package that now requires justification: the compact convex stationary state body, the affine observer-data and shell-response readouts, a nonnegative stationary potential with an exact zero-hidden stationary locus, the common zero-locus Hessian that vanishes on zero directions and is coercively positive on hidden directions, the induced zero/hidden splitting and hidden body, the symmetry and time-reversal structure, the exact zero-response shell realization on the zero locus, fixed-data shell solvability on that locus, the inert zero-hidden kernel and projector, and the stationary-Hessian coercivity inequality. Within this reduction class those stationary constituents are not merely available. They are forced, and the induced stationary law is unique up to the harmless orthogonal reparameterization of the hidden reduction target that leaves the stationary potential and Hessian unchanged. The benchmark package remains fixed throughout: the one operative metric is still exactly \(\CompletedMetric\), the ordinary matter route is unchanged, there is no second operative metric, and the low-energy matter law is not enlarged.
\end{abstract}

\section{Introduction}

The active derivational program of \TheoryName{} has again narrowed what counts as an honest next step. The primitive-observer-reduction datum theorem, the defect-fiber factorization theorem, the one-metric output theorem, the chain-forcing theorem, the selector theorem, the stationary-construction theorem, the explicit-package theorem, the primitive-normal-form theorem, the ambient-shell-law theorem, the combined-response theorem, the graph-collapse theorem, the completion-core theorem, the ambient-dynamics theorem, the selector-law theorem, and the later spectral-law theorem each discharged what was honestly open at its own layer. The latest result matters because it removes another tempting stopping point. Once the spectral package is itself a canonical reduction of a deeper stationary-Hessian law, it is no longer enough to revisit the spectral layer under a different presentation.

The remaining honest burden is therefore one layer deeper again. The task is now to explain why the zero-hidden stationary Hessian law should exist, rather than letting the compact convex stationary state body, the affine observer-data and shell-response readouts, the nonnegative stationary potential, the exact zero-hidden stationary locus, the common zero-locus Hessian, the zero/hidden splitting, the hidden body, the shell realization, the solvability statement, the inert kernel and projector, and the stationary-Hessian coercivity inequality stand as the new deepest disciplined assumptions. That is the burden addressed here.

The present note makes the smallest honest collapse that can presently be stated on the active line. It does not claim a final microscopic substrate law. Its gain is narrower and more upstream. It replaces an independently postulated stationary potential and common Hessian by a smaller stationary reduction law: one smooth reduction map whose exact zero locus and common zero-locus linearization canonically generate both the stationary potential and the stationary-Hessian coercivity. In that precise sense, the stationary-Hessian layer itself ceases to be an independent stopping point.

\paragraph{Framework and claim boundary.}
\TheoryProgram{} denotes the broader \Aether{} / \Aflow{} framework built on the claims that the \Aether{} is the underlying four-dimensional substrate of reality, the \Aflow{} is its intrinsic ordered motion, observed three-dimensional space is the local experiential slice of that deeper substrate, S-time is the experienced order of change arising from matter, light, and the \Aflow{}, and observed expansion is the three-dimensional appearance of deeper four-dimensional ordered motion. Within that framework, \TheoryName{} denotes the exact-closure theory in which the effective gravitational dynamics are adopted to be exactly Einsteinian with universal matter coupling. In the active sequence, \emph{adoption} means use of that established relativistic dynamics without claiming substrate derivation, while \emph{derivation} is reserved for a first-principles recovery from explicit substrate variables.

The present note stays strictly inside that boundary. It does not claim that final substrate microphysics has already been found. Its narrower theorem is only that, on the active primitive-observer-reduction line, the zero-hidden stationary Hessian law is itself forced by a smaller zero-hidden stationary reduction law, and is unique there modulo the harmless orthogonal reparameterization of the hidden reduction target that leaves the induced stationary potential and Hessian unchanged.

\section{Fixed Primitive Continuation Data}

\begin{definition}[Recorded primitive continuation data]
\label{def:fixed}
Fix the following continuation data from the active primitive-reservoir observer-reduction line.
\begin{enumerate}
    \item For each sector label \(\sigma\in\{\Car,\Cur,\Hom\}\), a primitive forcing shell
    \[
    \ForSpace{\sigma},
    \]
    a visible reduction map
    \[
    \Reduction{\sigma}:\ForSpace{\sigma}\to X_{\sigma},
    \]
    and shell-level current and equilibrium carriers
    \[
    \CurrentCarrier{\sigma}:\ForSpace{\sigma}\to Z_{\sigma}^{\mathrm{cur}},
    \qquad
    \EqCarrier{\sigma}:\ForSpace{\sigma}\to Z_{\sigma}^{\mathrm{eq}}.
    \]
    \item The product shell
    \[
    \PrimShell
    \equiv
    \ForSpace{\Car}\times \ForSpace{\Cur}\times \ForSpace{\Hom}.
    \]
    \item Fixed shell extractors
    \[
    \CurrentExtract{\sigma}:Z_{\sigma}^{\mathrm{cur}}\to Y_{\sigma}^{\mathrm{cur}},
    \qquad
    \EqExtract{\sigma}:Z_{\sigma}^{\mathrm{eq}}\to Y_{\sigma}^{\mathrm{eq}},
    \]
    together with the recorded shell composites
    \begin{equation}
    \PrimCurrent{\sigma}
    =
    \CurrentExtract{\sigma}\circ \CurrentCarrier{\sigma},
    \qquad
    \PrimTheta{\sigma}
    =
    \EqExtract{\sigma}\circ \EqCarrier{\sigma}.
    \label{eq:primcomposites}
    \end{equation}
    \item The recorded metric and ordinary-matter outputs
    \[
    \MetricOut:\PrimShell\to\{\text{observer metrics}\},
    \qquad
    \MatterOut:\PrimShell\times\Psi\to\{\text{observer stress tensors}\},
    \]
    satisfying
    \begin{equation}
    \MetricOut(\mathbf w)=\CompletedMetric,
    \qquad
    \MatterOut(\mathbf w,\psi)
    =
    T_{\mu\nu}^{\mathrm{matter}}[\CompletedMetric,\psi]
    \label{eq:fixedoutputs}
    \end{equation}
    for every \(\mathbf w\in\PrimShell\) and every matter configuration \(\psi\in\Psi\).
\end{enumerate}
\end{definition}

\begin{remark}
Definition \ref{def:fixed} is not reopened here. The primitive shell data and the benchmark one-metric / ordinary-matter outputs remain fixed continuation data. The work begins one layer deeper again: why the stationary-Hessian law that now carries the active line should itself be forced.
\end{remark}

\section{Target Stationary Hessian Law}

\begin{definition}[Zero-hidden stationary Hessian law]
\label{def:stationary}
Fix \(\sigma\in\{\Car,\Cur,\Hom\}\). A \emph{zero-hidden stationary Hessian law} for sector \(\sigma\) consists of the following data.
\begin{enumerate}
    \item A finite-dimensional Euclidean stationary state space
    \[
    \StatSpace{\sigma},
    \]
    the observer-data space
    \[
    \DataSpace{\sigma}
    \equiv
    \VisibleAffine{\sigma}\times Z_{\sigma}^{\mathrm{cur}}\times Z_{\sigma}^{\mathrm{eq}},
    \]
    an ambient shell-response space
    \[
    \AmbientTarget{\sigma},
    \]
    and a compact convex stationary state body
    \[
    \StatBody{\sigma}\subset\StatSpace{\sigma}
    \]
    with nonempty interior.
    \item An affine observer-data readout
    \[
    \StatData{\sigma}:\StatSpace{\sigma}\to\DataSpace{\sigma}
    \]
    and an affine shell-response readout
    \[
    \StatShell{\sigma}:\StatSpace{\sigma}\to\AmbientTarget{\sigma}.
    \]
    \item A smooth nonnegative stationary potential
    \[
    \StatPotential{\sigma}:\StatSpace{\sigma}\to[0,\infty)
    \]
    whose exact zero-hidden stationary locus
    \[
    \StatZero{\sigma}
    \equiv
    \left\{
    u\in\StatSpace{\sigma}:
    \StatPotential{\sigma}(u)=0
    \right\}
    \]
    is a linear subspace. Define the hidden stationary space
    \[
    \StatHidden{\sigma}
    \equiv
    \bigl(\StatZero{\sigma}\bigr)^{\perp},
    \]
    and let
    \[
    \StatProj{\sigma}:\StatSpace{\sigma}\to\StatHidden{\sigma}
    \]
    denote the orthogonal projector.
    \item A compact convex hidden stationary body
    \[
    \StatHiddenBody{\sigma}\subset\StatHidden{\sigma}
    \]
    with
    \[
    0\in\operatorname{int}\StatHiddenBody{\sigma},
    \]
    such that the stationary state body splits exactly as
    \begin{equation}
    \StatBody{\sigma}
    =
    \left\{
    z+\eta:
    z\in\StatBody{\sigma}\cap\StatZero{\sigma},
    \quad
    \eta\in\StatHiddenBody{\sigma}
    \right\}.
    \label{eq:statsplitting}
    \end{equation}
    \item A symmetric bilinear form
    \[
    \StatHess{\sigma}:
    \StatSpace{\sigma}\times\StatSpace{\sigma}\to\mathbb R
    \]
    such that for every \(z\in\StatZero{\sigma}\),
    \begin{equation}
    D\StatPotential{\sigma}(z)=0,
    \qquad
    D^{2}\StatPotential{\sigma}(z)[u,v]
    =
    \StatHess{\sigma}(u,v)
    \label{eq:commonhessian}
    \end{equation}
    for all \(u,v\in\StatSpace{\sigma}\). Assume moreover that
    \begin{equation}
    \StatHess{\sigma}(\delta z,u)=0
    \qquad
    \text{for all }\delta z\in\StatZero{\sigma},\;u\in\StatSpace{\sigma},
    \label{eq:hesszero}
    \end{equation}
    and that there exists \(\lambda_{\sigma}^{\mathrm{st}}>0\) with
    \begin{equation}
    \StatHess{\sigma}(\eta,\eta)
    \ge
    \lambda_{\sigma}^{\mathrm{st}}\,
    \|\eta\|^{2}
    \qquad
    \text{for all }\eta\in\StatHidden{\sigma}.
    \label{eq:stgap}
    \end{equation}
    \item A compact symmetry group
    \[
    \StatSym{\sigma}\subset
    O(\StatSpace{\sigma})\times
    O(\DataSpace{\sigma})\times
    O(\AmbientTarget{\sigma})
    \]
    and an orthogonal involution
    \[
    \StatTime{\sigma}\in
    O(\StatSpace{\sigma})\times
    O(\DataSpace{\sigma})\times
    O(\AmbientTarget{\sigma})
    \]
    preserving \(\StatBody{\sigma}\), \(\StatData{\sigma}\), \(\StatShell{\sigma}\), and \(\StatPotential{\sigma}\).
    \item A smooth embedding
    \[
    \jmath_{\sigma}:X_{\sigma}\hookrightarrow\VisibleAffine{\sigma}
    \]
    and a smooth zero-response shell realization
    \[
    \StatEmbed{\sigma}:\ForSpace{\sigma}\hookrightarrow\StatSpace{\sigma}
    \]
    whose image is exactly
    \[
    \StatBody{\sigma}\cap
    \StatZero{\sigma}\cap
    \bigl(\StatShell{\sigma}\bigr)^{-1}(0),
    \]
    and such that
    \begin{equation}
    \StatData{\sigma}\bigl(\StatEmbed{\sigma}(w)\bigr)
    =
    \bigl(
    \jmath_{\sigma}(\Reduction{\sigma}(w)),
    \CurrentCarrier{\sigma}(w),
    \EqCarrier{\sigma}(w)
    \bigr)
    \label{eq:statcompat}
    \end{equation}
    for every \(w\in\ForSpace{\sigma}\).
    \item Fixed-data shell solvability on the zero-hidden stationary locus:
    \begin{equation}
    \StatData{\sigma}\bigl(
    \StatBody{\sigma}\cap\StatZero{\sigma}
    \bigr)
    =
    \StatData{\sigma}\Bigl(
    \StatBody{\sigma}\cap
    \StatZero{\sigma}\cap
    \bigl(\StatShell{\sigma}\bigr)^{-1}(0)
    \Bigr).
    \label{eq:statsolvability}
    \end{equation}
    \item The inert zero-hidden kernel
    \[
    \Kernel{\sigma}
    \equiv
    \ker\bigl(D\StatData{\sigma}\big|_{\StatZero{\sigma}}\bigr)
    \cap
    \ker\bigl(D\StatShell{\sigma}\big|_{\StatZero{\sigma}}\bigr)
    \]
    and the orthogonal projector
    \[
    \StateComp{\sigma}:
    \StatZero{\sigma}\to
    {\Kernel{\sigma}}^{\perp}\cap\StatZero{\sigma}.
    \]
    \item Stationary-Hessian coercivity on data-invisible directions: there exists \(\kappa_{\sigma}^{\mathrm{st}}>0\) such that for every
    \[
    w\in\ForSpace{\sigma},
    \qquad
    \delta u\in T_{\StatEmbed{\sigma}(w)}\StatSpace{\sigma}=\StatSpace{\sigma}
    \]
    satisfying
    \[
    D\StatData{\sigma}\,\delta u=0,
    \]
    one has
    \begin{equation}
    \|D\StatShell{\sigma}\,\delta u\|^{2}
    +
    \StatHess{\sigma}(\delta u,\delta u)
    \ge
    \kappa_{\sigma}^{\mathrm{st}}
    \left(
    \|\StateComp{\sigma}\bigl((I-\StatProj{\sigma})\delta u\bigr)\|^{2}
    +
    \|\StatProj{\sigma}\delta u\|^{2}
    \right).
    \label{eq:statcoercive}
    \end{equation}
\end{enumerate}
\end{definition}

\begin{remark}
Definition \ref{def:stationary} records exactly the stationary-layer objects that now require deeper justification. The present note does not enlarge that package. It asks whether those same objects can be forced by a smaller stationary reduction law beneath them.
\end{remark}

\section{Zero-Hidden Stationary Reduction Law}

\begin{definition}[Zero-hidden stationary reduction law]
\label{def:reduction}
Fix \(\sigma\in\{\Car,\Cur,\Hom\}\). A \emph{zero-hidden stationary reduction law} for sector \(\sigma\) consists of the following data.
\begin{enumerate}
    \item A finite-dimensional Euclidean reduction state space
    \[
    \RedSpace{\sigma},
    \]
    the observer-data space
    \[
    \DataSpace{\sigma}
    \equiv
    \VisibleAffine{\sigma}\times Z_{\sigma}^{\mathrm{cur}}\times Z_{\sigma}^{\mathrm{eq}},
    \]
    an ambient shell-response space
    \[
    \AmbientTarget{\sigma},
    \]
    a finite-dimensional Euclidean hidden reduction target
    \[
    \RedTarget{\sigma},
    \]
    and a compact convex reduction state body
    \[
    \RedBody{\sigma}\subset\RedSpace{\sigma}
    \]
    with nonempty interior.
    \item An affine observer-data readout
    \[
    \RedData{\sigma}:\RedSpace{\sigma}\to\DataSpace{\sigma}
    \]
    and an affine shell-response readout
    \[
    \RedShell{\sigma}:\RedSpace{\sigma}\to\AmbientTarget{\sigma}.
    \]
    \item A smooth stationary reduction map
    \[
    \RedMap{\sigma}:\RedSpace{\sigma}\to\RedTarget{\sigma}
    \]
    whose exact zero-hidden reduction locus
    \[
    \RedZero{\sigma}
    \equiv
    \left\{
    u\in\RedSpace{\sigma}:
    \RedMap{\sigma}(u)=0
    \right\}
    \]
    is a linear subspace. Define the hidden reduction space
    \[
    \RedHidden{\sigma}
    \equiv
    \bigl(\RedZero{\sigma}\bigr)^{\perp},
    \]
    and let
    \[
    \RedProj{\sigma}:\RedSpace{\sigma}\to\RedHidden{\sigma}
    \]
    denote the orthogonal projector.
    \item A compact convex hidden reduction body
    \[
    \RedHiddenBody{\sigma}\subset\RedHidden{\sigma}
    \]
    with
    \[
    0\in\operatorname{int}\RedHiddenBody{\sigma},
    \]
    such that the reduction state body splits exactly as
    \begin{equation}
    \RedBody{\sigma}
    =
    \left\{
    z+\eta:
    z\in\RedBody{\sigma}\cap\RedZero{\sigma},
    \quad
    \eta\in\RedHiddenBody{\sigma}
    \right\}.
    \label{eq:redsplitting}
    \end{equation}
    \item A common zero-locus linearization
    \[
    \RedLin{\sigma}:\RedSpace{\sigma}\to\RedTarget{\sigma}
    \]
    such that for every \(z\in\RedZero{\sigma}\),
    \begin{equation}
    D\RedMap{\sigma}(z)=\RedLin{\sigma}.
    \label{eq:redlinearization}
    \end{equation}
    Assume moreover that
    \begin{equation}
    \RedLin{\sigma}(\delta z)=0
    \qquad
    \text{for all }\delta z\in\RedZero{\sigma},
    \label{eq:redzero}
    \end{equation}
    and that there exists \(\lambda_{\sigma}^{\mathrm{red}}>0\) with
    \begin{equation}
    \|\RedLin{\sigma}\eta\|_{\RedTarget{\sigma}}^{2}
    \ge
    \lambda_{\sigma}^{\mathrm{red}}\,
    \|\eta\|^{2}
    \qquad
    \text{for all }\eta\in\RedHidden{\sigma}.
    \label{eq:redgap}
    \end{equation}
    \item A compact symmetry group
    \[
    \RedSym{\sigma}\subset
    O(\RedSpace{\sigma})\times
    O(\DataSpace{\sigma})\times
    O(\AmbientTarget{\sigma})
    \]
    and an orthogonal involution
    \[
    \RedTime{\sigma}\in
    O(\RedSpace{\sigma})\times
    O(\DataSpace{\sigma})\times
    O(\AmbientTarget{\sigma})
    \]
    preserving \(\RedBody{\sigma}\), \(\RedData{\sigma}\), \(\RedShell{\sigma}\), and \(\RedMap{\sigma}\).
    \item A smooth embedding
    \[
    \jmath_{\sigma}:X_{\sigma}\hookrightarrow\VisibleAffine{\sigma}
    \]
    and a smooth zero-response shell realization
    \[
    \RedEmbed{\sigma}:\ForSpace{\sigma}\hookrightarrow\RedSpace{\sigma}
    \]
    whose image is exactly
    \[
    \RedBody{\sigma}\cap
    \RedZero{\sigma}\cap
    \bigl(\RedShell{\sigma}\bigr)^{-1}(0),
    \]
    and such that
    \begin{equation}
    \RedData{\sigma}\bigl(\RedEmbed{\sigma}(w)\bigr)
    =
    \bigl(
    \jmath_{\sigma}(\Reduction{\sigma}(w)),
    \CurrentCarrier{\sigma}(w),
    \EqCarrier{\sigma}(w)
    \bigr)
    \label{eq:redcompat}
    \end{equation}
    for every \(w\in\ForSpace{\sigma}\).
    \item Fixed-data shell solvability on the zero-hidden reduction locus:
    \begin{equation}
    \RedData{\sigma}\bigl(
    \RedBody{\sigma}\cap\RedZero{\sigma}
    \bigr)
    =
    \RedData{\sigma}\Bigl(
    \RedBody{\sigma}\cap
    \RedZero{\sigma}\cap
    \bigl(\RedShell{\sigma}\bigr)^{-1}(0)
    \Bigr).
    \label{eq:redsolvability}
    \end{equation}
    \item The inert zero-hidden kernel
    \[
    \Kernel{\sigma}
    \equiv
    \ker\bigl(D\RedData{\sigma}\big|_{\RedZero{\sigma}}\bigr)
    \cap
    \ker\bigl(D\RedShell{\sigma}\big|_{\RedZero{\sigma}}\bigr)
    \]
    and the orthogonal projector
    \[
    \StateComp{\sigma}:
    \RedZero{\sigma}\to
    {\Kernel{\sigma}}^{\perp}\cap\RedZero{\sigma}.
    \]
    \item Reduction coercivity on data-invisible directions: there exists \(\kappa_{\sigma}^{\mathrm{red}}>0\) such that for every
    \[
    w\in\ForSpace{\sigma},
    \qquad
    \delta u\in T_{\RedEmbed{\sigma}(w)}\RedSpace{\sigma}=\RedSpace{\sigma}
    \]
    satisfying
    \[
    D\RedData{\sigma}\,\delta u=0,
    \]
    one has
    \begin{equation}
    \|D\RedShell{\sigma}\,\delta u\|^{2}
    +
    \|\RedLin{\sigma}\delta u\|_{\RedTarget{\sigma}}^{2}
    \ge
    \kappa_{\sigma}^{\mathrm{red}}
    \left(
    \|\StateComp{\sigma}\bigl((I-\RedProj{\sigma})\delta u\bigr)\|^{2}
    +
    \|\RedProj{\sigma}\delta u\|^{2}
    \right).
    \label{eq:redcoercive}
    \end{equation}
\end{enumerate}
\end{definition}

\begin{remark}
Definition \ref{def:reduction} is smaller than Definition \ref{def:stationary} in exactly the way now needed. The nonnegative stationary potential is not stipulated independently, and neither is the common stationary Hessian. Both are extracted from one smooth reduction map together with its common zero-locus linearization.
\end{remark}

\section{Canonical Stationary Law Induced by the Stationary Reduction}

\begin{definition}[Canonical stationary law induced by the stationary reduction]
\label{def:canonical}
Assume Definition \ref{def:reduction}. Define the canonical zero-hidden stationary Hessian law by:
\begin{enumerate}
    \item
    \[
    \StatSpace{\sigma}\equiv\RedSpace{\sigma},
    \qquad
    \StatBody{\sigma}\equiv\RedBody{\sigma},
    \qquad
    \StatData{\sigma}\equiv\RedData{\sigma},
    \qquad
    \StatShell{\sigma}\equiv\RedShell{\sigma}.
    \]
    \item The stationary potential is the norm-square reduction potential
    \begin{equation}
    \StatPotential{\sigma}(u)
    \equiv
    \frac{1}{2}
    \|\RedMap{\sigma}(u)\|_{\RedTarget{\sigma}}^{2}.
    \label{eq:statpotential}
    \end{equation}
    \item
    \[
    \StatZero{\sigma}\equiv\RedZero{\sigma},
    \qquad
    \StatHidden{\sigma}\equiv\RedHidden{\sigma},
    \qquad
    \StatProj{\sigma}\equiv\RedProj{\sigma},
    \qquad
    \StatHiddenBody{\sigma}\equiv\RedHiddenBody{\sigma}.
    \]
    \item The common stationary Hessian is the Gram form of the common zero-locus linearization:
    \begin{equation}
    \StatHess{\sigma}(u,v)
    \equiv
    \ev{\RedLin{\sigma}u,\RedLin{\sigma}v}_{\RedTarget{\sigma}}.
    \label{eq:gram}
    \end{equation}
    \item
    \[
    \StatSym{\sigma}\equiv\RedSym{\sigma},
    \qquad
    \StatTime{\sigma}\equiv\RedTime{\sigma},
    \qquad
    \StatEmbed{\sigma}\equiv\RedEmbed{\sigma}.
    \]
\end{enumerate}
\end{definition}

\begin{lemma}[The stationary reduction law determines the stationary potential and Hessian]
\label{lem:potential}
Assume Definition \ref{def:reduction}. Then the canonical data of Definition \ref{def:canonical} satisfy:
\begin{enumerate}
    \item \(\StatPotential{\sigma}\ge 0\) on \(\StatSpace{\sigma}\), and
    \[
    \StatZero{\sigma}
    =
    \left\{
    u\in\StatSpace{\sigma}:
    \StatPotential{\sigma}(u)=0
    \right\}
    =
    \RedZero{\sigma};
    \]
    \item for every \(z\in\StatZero{\sigma}\),
    \[
    D\StatPotential{\sigma}(z)=0,
    \qquad
    D^{2}\StatPotential{\sigma}(z)[u,v]
    =
    \StatHess{\sigma}(u,v)
    \quad
    \text{for all }u,v\in\StatSpace{\sigma};
    \]
    \item
    \[
    \StatHess{\sigma}(\delta z,u)=0
    \qquad
    \text{for all }\delta z\in\StatZero{\sigma},\;u\in\StatSpace{\sigma};
    \]
    \item
    \[
    \StatHess{\sigma}(\eta,\eta)
    =
    \|\RedLin{\sigma}\eta\|_{\RedTarget{\sigma}}^{2}
    \ge
    \lambda_{\sigma}^{\mathrm{red}}\|\eta\|^{2}
    \qquad
    \text{for all }\eta\in\StatHidden{\sigma}.
    \]
\end{enumerate}
\end{lemma}

\begin{proof}
Item (1) is immediate from \eqref{eq:statpotential}. Moreover,
\[
\StatPotential{\sigma}(u)=0
\iff
\RedMap{\sigma}(u)=0
\iff
u\in\RedZero{\sigma},
\]
so the exact zero-hidden stationary locus is precisely \(\RedZero{\sigma}\).

For item (2), differentiate \eqref{eq:statpotential}:
\[
D\StatPotential{\sigma}(u)[\delta u]
=
\ev{\RedMap{\sigma}(u),D\RedMap{\sigma}(u)\delta u}_{\RedTarget{\sigma}}.
\]
If \(z\in\RedZero{\sigma}\), then \(\RedMap{\sigma}(z)=0\), hence \(D\StatPotential{\sigma}(z)=0\). Differentiating once more and again using \(\RedMap{\sigma}(z)=0\) gives
\[
D^{2}\StatPotential{\sigma}(z)[u,v]
=
\ev{D\RedMap{\sigma}(z)u,D\RedMap{\sigma}(z)v}_{\RedTarget{\sigma}}
=
\ev{\RedLin{\sigma}u,\RedLin{\sigma}v}_{\RedTarget{\sigma}}
=
\StatHess{\sigma}(u,v),
\]
where \eqref{eq:redlinearization} is used in the second equality and \eqref{eq:gram} in the third.

Item (3) follows from \eqref{eq:redzero}:
\[
\StatHess{\sigma}(\delta z,u)
=
\ev{\RedLin{\sigma}\delta z,\RedLin{\sigma}u}_{\RedTarget{\sigma}}
=
0.
\]
Item (4) is exactly \eqref{eq:redgap} rewritten using \eqref{eq:gram} and the identification \(\StatHidden{\sigma}=\RedHidden{\sigma}\).
\end{proof}

\begin{proposition}[The stationary reduction law induces a zero-hidden stationary Hessian law]
\label{prop:stationary}
Assume Definitions \ref{def:reduction} and \ref{def:canonical}. Then the canonical data of Definition \ref{def:canonical} satisfy every requirement of Definition \ref{def:stationary}.
\end{proposition}

\begin{proof}
Items (1) and (2) of Definition \ref{def:stationary} are immediate from Definition \ref{def:canonical}. By Lemma \ref{lem:potential}, \(\StatPotential{\sigma}\) is nonnegative and has exact zero-hidden stationary locus \(\StatZero{\sigma}=\RedZero{\sigma}\), so item (3) holds with
\[
\StatHidden{\sigma}
=
\bigl(\StatZero{\sigma}\bigr)^{\perp}
=
\RedHidden{\sigma},
\qquad
\StatProj{\sigma}=\RedProj{\sigma}.
\]

The compact convex hidden stationary body is \(\StatHiddenBody{\sigma}=\RedHiddenBody{\sigma}\), and the stationary body splitting \eqref{eq:statsplitting} is exactly the reduction splitting \eqref{eq:redsplitting}. This gives item (4). Item (5) is precisely Lemma \ref{lem:potential}.

For item (6), every element of \(\RedSym{\sigma}\) and \(\RedTime{\sigma}\) preserves \(\RedMap{\sigma}\), so by \eqref{eq:statpotential} it preserves \(\StatPotential{\sigma}\) as well. Thus the symmetry and time-reversal data preserve the full stationary law.

The shell realization and data compatibility in item (7) are exactly \eqref{eq:redcompat}, because
\[
\StatBody{\sigma}\cap\StatZero{\sigma}\cap\bigl(\StatShell{\sigma}\bigr)^{-1}(0)
=
\RedBody{\sigma}\cap\RedZero{\sigma}\cap\bigl(\RedShell{\sigma}\bigr)^{-1}(0).
\]
Fixed-data shell solvability \eqref{eq:statsolvability} is exactly \eqref{eq:redsolvability}. The inert zero-hidden kernel and orthogonal projector in item (9) are defined by the same formulas, since
\[
\StatData{\sigma}=\RedData{\sigma},
\qquad
\StatShell{\sigma}=\RedShell{\sigma},
\qquad
\StatZero{\sigma}=\RedZero{\sigma}.
\]

For item (10), let
\[
w\in\ForSpace{\sigma},
\qquad
\delta u\in T_{\StatEmbed{\sigma}(w)}\StatSpace{\sigma}=\RedSpace{\sigma}
\]
with \(D\StatData{\sigma}\,\delta u=0\). Then
\[
\StatHess{\sigma}(\delta u,\delta u)
=
\|\RedLin{\sigma}\delta u\|_{\RedTarget{\sigma}}^{2}
\]
by \eqref{eq:gram}. Applying \eqref{eq:redcoercive} therefore yields \eqref{eq:statcoercive} with \(\kappa_{\sigma}^{\mathrm{st}}=\kappa_{\sigma}^{\mathrm{red}}\). All requirements of Definition \ref{def:stationary} are satisfied.
\end{proof}

\section{Forced Constituents and Uniqueness}

\begin{proposition}[All stationary constituents are forced by the stationary reduction law]
\label{prop:forced}
Assume Definition \ref{def:reduction}. Then the following stationary-level constituents are fixed by the stationary reduction law itself.
\begin{enumerate}
    \item The stationary state body is necessarily the reduction state body \(\StatBody{\sigma}=\RedBody{\sigma}\).
    \item The affine observer-data and shell-response readouts are necessarily \(\StatData{\sigma}=\RedData{\sigma}\) and \(\StatShell{\sigma}=\RedShell{\sigma}\).
    \item The stationary potential is necessarily the norm-square reduction potential
    \[
    \StatPotential{\sigma}(u)
    =
    \frac{1}{2}\|\RedMap{\sigma}(u)\|_{\RedTarget{\sigma}}^{2},
    \]
    so the exact zero-hidden stationary locus is necessarily
    \[
    \StatZero{\sigma}=\RedZero{\sigma}.
    \]
    \item The common zero-locus stationary Hessian is necessarily the Gram form of the common zero-locus linearization:
    \[
    \StatHess{\sigma}(u,v)
    =
    \ev{\RedLin{\sigma}u,\RedLin{\sigma}v}_{\RedTarget{\sigma}}.
    \]
    Consequently it vanishes on zero directions and is coercively positive on hidden directions.
    \item The zero/hidden splitting is necessarily the reduction splitting
    \[
    \StatSpace{\sigma}
    =
    \StatZero{\sigma}\oplus\StatHidden{\sigma}
    =
    \RedZero{\sigma}\oplus\RedHidden{\sigma},
    \]
    the hidden body is necessarily \(\StatHiddenBody{\sigma}=\RedHiddenBody{\sigma}\), and the orthogonal projector onto the hidden space is necessarily \(\StatProj{\sigma}=\RedProj{\sigma}\).
    \item The symmetry and time-reversal structure, the exact zero-response shell realization, fixed-data shell solvability on the zero locus, the inert zero-hidden kernel and projector, and the stationary-Hessian coercivity inequality are all induced by the same reduction law and are not additional stationary freedom.
\end{enumerate}
\end{proposition}

\begin{proof}
Items (1) and (2) are immediate from Definition \ref{def:canonical}. Item (3) is \eqref{eq:statpotential}. Item (4) is Lemma \ref{lem:potential}: once the common zero-locus linearization \(\RedLin{\sigma}\) is fixed, the stationary Hessian is uniquely the associated Gram form \eqref{eq:gram}, and its vanishing on zero directions and hidden-sector coercivity are forced. Item (5) is part of Definition \ref{def:canonical}. Item (6) is exactly the content of Proposition \ref{prop:stationary}: every remaining stationary constituent descends directly from the same stationary-reduction package.
\end{proof}

\begin{theorem}[Zero-hidden stationary Hessian law necessity / uniqueness from stationary reduction law]
\label{thm:necessity}
Assume Definition \ref{def:reduction}. Then:
\begin{enumerate}
    \item the canonical construction of Definition \ref{def:canonical} yields a zero-hidden stationary Hessian law in the sense of Definition \ref{def:stationary};
    \item every stationary-level object singled out by the previous spectral-law theorem is forced by the reduction law: the compact convex stationary state body, the affine observer-data and shell-response readouts, the nonnegative stationary potential with exact zero-hidden stationary locus, the common zero-locus Hessian that vanishes on zero directions and is coercively positive on hidden directions, the induced zero/hidden splitting and hidden body, the symmetry and time-reversal structure, the exact zero-response shell realization on the zero locus, fixed-data shell solvability on that locus, the inert zero-hidden kernel and projector, and the stationary-Hessian coercivity inequality;
    \item any orthogonal reparameterization of the hidden reduction target \(\RedTarget{\sigma}\) leaves the induced stationary law unchanged, and there is no further stationary-Hessian freedom once the reduction law is fixed;
    \item consequently the stationary-Hessian layer is not an independent remaining freedom once the stationary-reduction law is fixed.
\end{enumerate}
\end{theorem}

\begin{proof}
Item (1) is Proposition \ref{prop:stationary}. Item (2) is Proposition \ref{prop:forced}. For item (3), if \(Q:\RedTarget{\sigma}\to\RedTarget{\sigma}\) is orthogonal and one replaces
\[
\RedMap{\sigma}\mapsto Q\RedMap{\sigma},
\qquad
\RedLin{\sigma}\mapsto Q\RedLin{\sigma},
\]
then \eqref{eq:statpotential} and \eqref{eq:gram} are unchanged because \(Q\) preserves the Euclidean norm and inner product on \(\RedTarget{\sigma}\). Conversely, the canonical formulas \eqref{eq:statpotential} and \eqref{eq:gram} show that once the reduction law is fixed, the stationary potential and the common stationary Hessian are fixed with no extra freedom. Item (4) is the routing consequence of the previous items.
\end{proof}

\section{Benchmark Preservation}

\begin{theorem}[The derived stationary law preserves the same one-metric benchmark package]
\label{thm:outputs}
Assume Definition \ref{def:reduction}. Then the benchmark output statements of Definition \ref{def:fixed} remain unchanged:
\[
\MetricOut(\mathbf w)=\CompletedMetric,
\qquad
\MatterOut(\mathbf w,\psi)
=
T_{\mu\nu}^{\mathrm{matter}}[\CompletedMetric,\psi]
\]
for every \(\mathbf w\in\PrimShell\) and every matter configuration \(\psi\in\Psi\). Consequently the deeper stationary-reduction law and the induced stationary-Hessian law introduce neither a second operative metric nor an enlarged low-energy matter law.
\end{theorem}

\begin{proof}
The present note changes only the upstream origin of the stationary data. It does not alter the primitive shell \(\PrimShell\), the recorded metric map \(\MetricOut\), or the recorded matter map \(\MatterOut\). Those remain the fixed continuation data of Definition \ref{def:fixed}, so \eqref{eq:fixedoutputs} continues to hold exactly.
\end{proof}

\section{Main Theorem}

\begin{theorem}[The remaining burden moves below the stationary-Hessian layer]
\label{thm:main}
Assume the fixed primitive continuation data of Definition \ref{def:fixed} and, for each sector \(\sigma\in\{\Car,\Cur,\Hom\}\), a zero-hidden stationary reduction law in the sense of Definition \ref{def:reduction}. Then:
\begin{enumerate}
    \item the zero-hidden stationary Hessian law of Definition \ref{def:stationary} is canonically induced by the stationary-reduction law;
    \item the compact convex stationary state body, the affine observer-data and shell-response readouts, the nonnegative stationary potential with exact zero-hidden stationary locus, the common zero-locus Hessian that vanishes on zero directions and is coercively positive on hidden directions, the induced zero/hidden splitting and hidden body, the exact zero-response shell realization, fixed-data shell solvability on the zero locus, the inert zero-hidden kernel and projector, and the stationary-Hessian coercivity inequality are all forced by that same stationary-reduction law;
    \item the one operative metric remains exactly \(\CompletedMetric\), the ordinary matter route remains exactly the standard one through that metric, there is no second operative metric, and the low-energy matter law is not enlarged;
    \item consequently the benchmark-facing burden after the previous spectral-law theorem is discharged in the stronger form now available on the active line: the zero-hidden stationary Hessian law is no longer an unexplained stopping point, but a necessary reduction of a smaller zero-hidden stationary reduction law.
\end{enumerate}
\end{theorem}

\begin{proof}
Item (1) is Theorem \ref{thm:necessity}(1). Item (2) is Theorem \ref{thm:necessity}(2). Item (3) is Theorem \ref{thm:outputs}. Item (4) is the combined routing consequence.
\end{proof}

\begin{corollary}[What remains open after this theorem]
\label{cor:next}
After Theorem \ref{thm:main}, the remaining honest burden is no longer to justify the zero-hidden stationary Hessian law itself on the active primitive-observer-reduction line. The next honest burden is deeper again: derive or justify the zero-hidden stationary reduction law from still more primitive substrate dynamics, or else prove a sharper inevitability theorem that removes even that remaining stationary-reduction-level freedom while preserving the same benchmark one-metric package and claim boundary.
\end{corollary}

\begin{proof}
Theorem \ref{thm:main} converts the stationary-Hessian layer from a remaining benchmark-facing burden into a canonical reduction of a deeper stationary-reduction law. What remains open is why that stationary-reduction law itself should hold, rather than becoming the new disciplined stopping point.
\end{proof}

\section{Conclusion}

The positive result of this note is deliberately narrower than a completed final microphysical derivation and stronger than another deeper same-output relay. The previous theorem had already shown that the spectral-law layer is a forced reduction of a stationary-Hessian law. The present note goes one honest step further by showing that the stationary-Hessian layer itself is not left hanging either. It is canonically induced by a smaller zero-hidden stationary reduction law.

That gain directly addresses the precise objects that had become the next burden: the compact convex stationary state body, the affine observer-data and shell-response readouts, the nonnegative stationary potential with exact zero-hidden stationary locus, the common zero-locus Hessian that vanishes on zero directions and is coercively positive on hidden directions, the induced zero/hidden splitting and hidden body, the symmetry and time-reversal structure, the exact zero-response shell realization on the zero locus, fixed-data shell solvability on that locus, the inert zero-hidden kernel and projector, and the stationary-Hessian coercivity inequality. Within the stationary-reduction class those are no longer optional inserted ingredients. They are forced outputs, and the only harmless freedom left is the orthogonal reparameterization of the hidden reduction target that does not change the induced stationary potential or Hessian.

The benchmark boundary remains fixed throughout. The same primitive shell remains the shell carrying the observer outputs, so the one operative metric is still exactly \(\CompletedMetric\), the ordinary matter route is unchanged, there is no second operative metric, and the low-energy matter law is not enlarged. This is therefore a disciplined upstream derivational gain, not a final completion claim. The note still does not derive the stationary-reduction law from ultimate substrate microphysics. That is the next honest open burden. But the stationary-Hessian layer is no longer the stopping point on the active line.

\end{document}
